% --- Archon physics formalization source begin ---
% archon:physics
% archon:covers PhyXMiniProblems/problem_phyx_mini_0121.lean
% archon:source-report reports/phyx_mini/problem_phyx_mini_0121.source.json

\chapter{Physics problem phyx\_mini\_0121}
\label{ch:PhyXMiniProblems_problem_phyx_mini_0121}

\paragraph{Problem source.}
\#\# Physical scenario

Figure shows an interferometer known as Fresnel’s biprism.The magnitude of the prism angle \$A\$ is extremely small. The green light with a wavelength of 500 nm is on a screen 2.00 m from the biprism.Take \$a = 0.200 \textbackslash{}, 	ext\{m\}\$, \$A = 3.50 \textbackslash{}, 	ext\{mrad\}\$, and \$n = 1.50\$.

\#\# Question

Calculate the spacing of the fringes of green light.

\#\# Answer choices

- A: A: \$1.57 \textbackslash{}times 10\textasciicircum{}\{-3\}\$ m

- B: B: \$1.57 \textbackslash{}times 10\textasciicircum{}\{-2\}\$ m

- C: C: \$1.57 \textbackslash{}times 10\textasciicircum{}\{-4\}\$ m

- D: D: \$1.57 \textbackslash{}times 10\textasciicircum{}\{-5\}\$ m

\#\# Figure caption (auxiliary; use the image as primary evidence)

The image depicts a diagram of a double-slit experiment with a lens. Here's a breakdown:

- **Light Sources and Slits**:
  - Three slits labeled \textbackslash{}( S\_1 \textbackslash{}), \textbackslash{}( S\_0 \textbackslash{}), and \textbackslash{}( S\_2 \textbackslash{}). 
  - These slits are aligned vertically with \textbackslash{}( S\_0 \textbackslash{}) in the center.

- **Lens**:
  - Positioned after the slits and marked as \textbackslash{}( A \textbackslash{}).
  - The lens refracts light rays emerging from the slits.

- **Light Rays**:
  - Purple lines with arrows indicating their direction.
  - Rays originate from the slits \textbackslash{}( S\_1 \textbackslash{}), \textbackslash{}( S\_0 \textbackslash{}), and \textbackslash{}( S\_2 \textbackslash{}) and pass through the lens \textbackslash{}( A \textbackslash{}).

- **Screen**:
  - On the right, perpendicular to the direction of the light rays.
  - Points \textbackslash{}( P \textbackslash{}) and \textbackslash{}( O \textbackslash{}) are marked where the rays converge after passing through the lens.

- **Distances**:
  - \textbackslash{}( d \textbackslash{}) is the vertical distance between the slits.
  - \textbackslash{}( a \textbackslash{}) is the horizontal distance from the slits to the lens.
  - \textbackslash{}( b \textbackslash{}) is the horizontal distance from the lens to the screen.

The arrangement illustrates the interference pattern on the screen formed by the light passing through the slits and focused by the lens.

\#\# Dataset metadata

Wave Optics; Physical Model Grounding Reasoning, Multi-Formula Reasoning

\paragraph{Recorded answer/context.}
A: A: \$1.57 \textbackslash{}times 10\textasciicircum{}\{-3\}\$ m

\paragraph{Figure/image path.}
/root/proposal\_for\_physic/science-mango/phyx\_mini\_run/phyx\_data/test\_image/121.png

\paragraph{Formalization target.}
create a compiling Lean file with sorry bodies at `PhyXMiniProblems/problem\_phyx\_mini\_0121.lean`. The Lean declarations must preserve the physical quantities, dimensions or dimensional roles, figure labels, governing-law hypotheses, and final relation expressed by this problem.
Use Mathlib/Physlib names found through LeanExplore where available. If a physics API is missing, introduce faithful local abstractions rather than scalar placeholder aliases.

\begin{theorem}[Physics formalization target]
\label{thm:physics:phyx_mini_0121:target}
\lean{PhyXMiniProblems.ProblemPhyXMini0121.fringeSpacing_matches_recordedAnswerA}
\uses{def:physics:phyx-mini-0121:phyxminiproblems-problemphyxmini0121-fresnelbiprisminterferometer, def:physics:phyx-mini-0121:phyxminiproblems-problemphyxmini0121-hasstatedproblemreadouts, def:physics:phyx-mini-0121:phyxminiproblems-problemphyxmini0121-hasphysicalparameters, def:physics:phyx-mini-0121:phyxminiproblems-problemphyxmini0121-satisfiescontrolledthinbiprismimagelaw, def:physics:phyx-mini-0121:phyxminiproblems-problemphyxmini0121-satisfiescontrolledadjacentfringelaw, def:physics:phyx-mini-0121:phyxminiproblems-problemphyxmini0121-matchesdisplayedprecision}
The assigned autoformalize agent should translate this physics problem into Lean declarations in the covered file, with theorem and lemma proof bodies written as `by sorry`.
\end{theorem}
\begin{proof}
This is an autoformalization task, not a proof task. Produce statements that can later be proved by the physics prover without weakening the physical model.
\end{proof}

% --- Archon declaration topology begin ---
\section{Declaration topology}
\begin{definition}[Dim Length]
  \label{def:physics:phyx-mini-0121:phyxminiproblems-problemphyxmini0121-dimlength}
  \lean{PhyXMiniProblems.ProblemPhyXMini0121.DimLength}
  A signed physical length represented independently of a chosen unit system
\end{definition}
\begin{proof}
  This is the declared model definition of Dim Length.
\end{proof}
\begin{definition}[length Value In]
  \label{def:physics:phyx-mini-0121:phyxminiproblems-problemphyxmini0121-lengthvaluein}
  \lean{PhyXMiniProblems.ProblemPhyXMini0121.lengthValueIn}
  \uses{def:physics:phyx-mini-0121:phyxminiproblems-problemphyxmini0121-dimlength}
  Read a physical length as a real scalar in the selected length unit
\end{definition}
\begin{proof}
  This is the declared model definition of length Value In.
\end{proof}
\begin{definition}[meters Value]
  \label{def:physics:phyx-mini-0121:phyxminiproblems-problemphyxmini0121-metersvalue}
  \lean{PhyXMiniProblems.ProblemPhyXMini0121.metersValue}
  \uses{def:physics:phyx-mini-0121:phyxminiproblems-problemphyxmini0121-dimlength, def:physics:phyx-mini-0121:phyxminiproblems-problemphyxmini0121-lengthvaluein}
  The metre readout used in the biprism geometry and interference laws
\end{definition}
\begin{proof}
  This is the declared model definition of meters Value.
\end{proof}
\begin{definition}[nanometers Value]
  \label{def:physics:phyx-mini-0121:phyxminiproblems-problemphyxmini0121-nanometersvalue}
  \lean{PhyXMiniProblems.ProblemPhyXMini0121.nanometersValue}
  \uses{def:physics:phyx-mini-0121:phyxminiproblems-problemphyxmini0121-dimlength, def:physics:phyx-mini-0121:phyxminiproblems-problemphyxmini0121-lengthvaluein}
  The nanometre readout used for the stated green-light wavelength
\end{definition}
\begin{proof}
  This is the declared model definition of nanometers Value.
\end{proof}
\begin{definition}[milliradians Value]
  \label{def:physics:phyx-mini-0121:phyxminiproblems-problemphyxmini0121-milliradiansvalue}
  \lean{PhyXMiniProblems.ProblemPhyXMini0121.milliradiansValue}
  The scalar milliradian readout of a physical angle
\end{definition}
\begin{proof}
  This is the declared model definition of milliradians Value.
\end{proof}
\begin{definition}[Source Label]
  \label{def:physics:phyx-mini-0121:phyxminiproblems-problemphyxmini0121-sourcelabel}
  \lean{PhyXMiniProblems.ProblemPhyXMini0121.SourceLabel}
  The three source labels in the primary figure
\end{definition}
\begin{proof}
  This is the declared model definition of Source Label.
\end{proof}
\begin{definition}[Source Role]
  \label{def:physics:phyx-mini-0121:phyxminiproblems-problemphyxmini0121-sourcerole}
  \lean{PhyXMiniProblems.ProblemPhyXMini0121.SourceRole}
  Whether a labelled source is actual or is a virtual biprism image
\end{definition}
\begin{proof}
  This is the declared model definition of Source Role.
\end{proof}
\begin{definition}[role]
  \label{def:physics:phyx-mini-0121:phyxminiproblems-problemphyxmini0121-sourcelabel-role}
  \lean{PhyXMiniProblems.ProblemPhyXMini0121.SourceLabel.role}
  \uses{def:physics:phyx-mini-0121:phyxminiproblems-problemphyxmini0121-sourcelabel, def:physics:phyx-mini-0121:phyxminiproblems-problemphyxmini0121-sourcerole}
  The optical role assigned to each source label by the ray diagram
\end{definition}
\begin{proof}
  This is the declared model definition of role.
\end{proof}
\begin{definition}[Screen Point]
  \label{def:physics:phyx-mini-0121:phyxminiproblems-problemphyxmini0121-screenpoint}
  \lean{PhyXMiniProblems.ProblemPhyXMini0121.ScreenPoint}
  The two explicitly marked points on the observation screen
\end{definition}
\begin{proof}
  This is the declared model definition of Screen Point.
\end{proof}
\begin{definition}[Optical Medium]
  \label{def:physics:phyx-mini-0121:phyxminiproblems-problemphyxmini0121-opticalmedium}
  \lean{PhyXMiniProblems.ProblemPhyXMini0121.OpticalMedium}
  The two homogeneous optical media relevant to the thin-prism deviation
\end{definition}
\begin{proof}
  This is the declared model definition of Optical Medium.
\end{proof}
\begin{definition}[Fresnel Biprism Interferometer]
  \label{def:physics:phyx-mini-0121:phyxminiproblems-problemphyxmini0121-fresnelbiprisminterferometer}
  \lean{PhyXMiniProblems.ProblemPhyXMini0121.FresnelBiprismInterferometer}
  \uses{def:physics:phyx-mini-0121:phyxminiproblems-problemphyxmini0121-dimlength, def:physics:phyx-mini-0121:phyxminiproblems-problemphyxmini0121-opticalmedium}
  The monochromatic source, Fresnel biprism, virtual images, and screen. sourceToBiprismDistance, biprismToScreenDistance, and virtualSourceSeparation are the figure labels a, b, and d. prismAngle is the common magnitude A displayed at both halves of the biprism. adjacentFringeSpacing is the physical quantity requested by the problem; it is stored independently and constrained only by the controlled approximation law below
\end{definition}
\begin{proof}
  This is the declared model definition of Fresnel Biprism Interferometer.
\end{proof}
\begin{definition}[Has Stated Problem Readouts]
  \label{def:physics:phyx-mini-0121:phyxminiproblems-problemphyxmini0121-hasstatedproblemreadouts}
  \lean{PhyXMiniProblems.ProblemPhyXMini0121.HasStatedProblemReadouts}
  \uses{def:physics:phyx-mini-0121:phyxminiproblems-problemphyxmini0121-metersvalue, def:physics:phyx-mini-0121:phyxminiproblems-problemphyxmini0121-nanometersvalue, def:physics:phyx-mini-0121:phyxminiproblems-problemphyxmini0121-milliradiansvalue, def:physics:phyx-mini-0121:phyxminiproblems-problemphyxmini0121-fresnelbiprisminterferometer}
  The numerical readouts given by the problem and primary figure: green light of wavelength 500 nm, a = 0.200 m, b = 2.00 m, common prism angle A = 3.50 mrad, air index one, and glass index n = 1.50
\end{definition}
\begin{proof}
  This is the declared model definition of Has Stated Problem Readouts.
\end{proof}
\begin{definition}[Has Physical Parameters]
  \label{def:physics:phyx-mini-0121:phyxminiproblems-problemphyxmini0121-hasphysicalparameters}
  \lean{PhyXMiniProblems.ProblemPhyXMini0121.HasPhysicalParameters}
  \uses{def:physics:phyx-mini-0121:phyxminiproblems-problemphyxmini0121-metersvalue, def:physics:phyx-mini-0121:phyxminiproblems-problemphyxmini0121-fresnelbiprisminterferometer}
  Positivity, optical-index ordering, and the small acute prism-angle branch
\end{definition}
\begin{proof}
  This is the declared model definition of Has Physical Parameters.
\end{proof}
\begin{definition}[Satisfies Controlled Thin Biprism Image Law]
  \label{def:physics:phyx-mini-0121:phyxminiproblems-problemphyxmini0121-satisfiescontrolledthinbiprismimagelaw}
  \lean{PhyXMiniProblems.ProblemPhyXMini0121.SatisfiesControlledThinBiprismImageLaw}
  \uses{def:physics:phyx-mini-0121:phyxminiproblems-problemphyxmini0121-metersvalue, def:physics:phyx-mini-0121:phyxminiproblems-problemphyxmini0121-fresnelbiprisminterferometer}
  A quantitative small-angle contract for the virtual-image separation. The leading term is the usual thin-prism prediction 2 a (n\_glass / n\_air - 1) A. Instead of globalizing that first-order formula as an exact optical law, this predicate exposes a metre-valued higher-order remainder and bounds it at cubic order in the dimensionless angle A. The bound is a modeling contract for the stated thin biprism, not the answer to the fringe-spacing question
\end{definition}
\begin{proof}
  This is the declared model definition of Satisfies Controlled Thin Biprism Image Law.
\end{proof}
\begin{definition}[Satisfies Controlled Adjacent Fringe Law]
  \label{def:physics:phyx-mini-0121:phyxminiproblems-problemphyxmini0121-satisfiescontrolledadjacentfringelaw}
  \lean{PhyXMiniProblems.ProblemPhyXMini0121.SatisfiesControlledAdjacentFringeLaw}
  \uses{def:physics:phyx-mini-0121:phyxminiproblems-problemphyxmini0121-metersvalue, def:physics:phyx-mini-0121:phyxminiproblems-problemphyxmini0121-fresnelbiprisminterferometer}
  A controlled near-axis interference contract for the coherent virtual sources S₁ and S₂. The leading term is λ (a + b) / d. The exact observed central adjacent spacing may differ by fringeRemainderMeters. Its relative error is bounded by the squared aperture ratio d / (2 (a + b)) plus the squared wavelength ratio λ / d, making the paraxial and many-fringe approximations explicit rather than asserting their leading term globally as an exact equality
\end{definition}
\begin{proof}
  This is the declared model definition of Satisfies Controlled Adjacent Fringe Law.
\end{proof}
\begin{definition}[Answer Choice]
  \label{def:physics:phyx-mini-0121:phyxminiproblems-problemphyxmini0121-answerchoice}
  \lean{PhyXMiniProblems.ProblemPhyXMini0121.AnswerChoice}
  Labels of the four multiple-choice answers printed with the problem
\end{definition}
\begin{proof}
  This is the declared model definition of Answer Choice.
\end{proof}
\begin{definition}[answer Fringe Spacing Meters]
  \label{def:physics:phyx-mini-0121:phyxminiproblems-problemphyxmini0121-answerfringespacingmeters}
  \lean{PhyXMiniProblems.ProblemPhyXMini0121.answerFringeSpacingMeters}
  \uses{def:physics:phyx-mini-0121:phyxminiproblems-problemphyxmini0121-answerchoice}
  Metre readout printed beside each displayed answer choice
\end{definition}
\begin{proof}
  This is the declared model definition of answer Fringe Spacing Meters.
\end{proof}
\begin{definition}[answer Rounding Tolerance Meters]
  \label{def:physics:phyx-mini-0121:phyxminiproblems-problemphyxmini0121-answerroundingtolerancemeters}
  \lean{PhyXMiniProblems.ProblemPhyXMini0121.answerRoundingToleranceMeters}
  \uses{def:physics:phyx-mini-0121:phyxminiproblems-problemphyxmini0121-answerchoice}
  Half a unit in the final displayed digit of each three-significant-figure choice
\end{definition}
\begin{proof}
  This is the declared model definition of answer Rounding Tolerance Meters.
\end{proof}
\begin{definition}[Matches Displayed Precision]
  \label{def:physics:phyx-mini-0121:phyxminiproblems-problemphyxmini0121-matchesdisplayedprecision}
  \lean{PhyXMiniProblems.ProblemPhyXMini0121.MatchesDisplayedPrecision}
  \uses{def:physics:phyx-mini-0121:phyxminiproblems-problemphyxmini0121-dimlength, def:physics:phyx-mini-0121:phyxminiproblems-problemphyxmini0121-metersvalue, def:physics:phyx-mini-0121:phyxminiproblems-problemphyxmini0121-answerchoice, def:physics:phyx-mini-0121:phyxminiproblems-problemphyxmini0121-answerfringespacingmeters, def:physics:phyx-mini-0121:phyxminiproblems-problemphyxmini0121-answerroundingtolerancemeters}
  A physical spacing agrees with a displayed choice to its printed precision
\end{definition}
\begin{proof}
  This is the declared model definition of Matches Displayed Precision.
\end{proof}
% --- Archon declaration topology end ---
% --- Archon physics formalization source end ---
